% ============================================================
% PathTriage — COMP9301 Midway Report
% Author: Hyeon (Tessa) Moon, z5660470
% Supervisor: Lachlan Jones
% Submitted: July 2026
% ============================================================
\documentclass[11pt,a4paper]{report}

% --- Packages ---
\usepackage[utf8]{inputenc}
\usepackage[T1]{fontenc}
\usepackage{mathptmx}              % Times-like serif font
\usepackage[margin=1in]{geometry}
\usepackage{graphicx}
\usepackage{xcolor}
\usepackage{hyperref}
\hypersetup{
  colorlinks=true,
  linkcolor=black,
  citecolor=black,
  urlcolor=blue
}
\usepackage{cite}                  % IEEE-style numeric citations [1], [2-4]
\usepackage{booktabs}              % nicer tables
\usepackage{listings}              % code blocks
\usepackage{caption}
\usepackage{setspace}
\onehalfspacing

% --- Metadata ---
\newcommand{\thesistitle}{PathTriage: Exploitability-Ranked IAM Attack-Path Discovery and Defender-Output Synthesis for AWS and Azure}
\newcommand{\authorname}{Hyeon (Tessa) Moon}
\newcommand{\studentid}{z5660470}
\newcommand{\supervisor}{Lachlan Jones}
\newcommand{\submissiondate}{Submitted: July 2026}

% ============================================================
\begin{document}

% --- Custom title page (matches the UNSW sample style) ---
\begin{titlepage}
\centering
\vspace*{2cm}

% \includegraphics[width=4cm]{unsw_logo.png}  % uncomment when logo file is added
\vspace{1.5cm}

\textbf{School of Computer Science and Engineering}\\[0.4cm]
\textbf{Faculty of Engineering}\\[0.4cm]
\textbf{The University of New South Wales}

\vspace{2.5cm}
{\Large \textbf{\thesistitle}}

\vspace{2.5cm}
by

\vspace{0.8cm}
{\Large \authorname}

\vspace{0.3cm}
\studentid

\vfill
\submissiondate \hspace{2em} Supervisor: \supervisor
\end{titlepage}

% ============================================================
% Preamble
% ============================================================
\chapter*{Preamble}
\addcontentsline{toc}{chapter}{Preamble}

\section*{Abstract}
This work investigates the construction of an exploitability-ranked catalogue of
IAM privilege-escalation paths in cloud environments, and the synthesis of
defender output --- detection queries and remediation policies --- for each ranked
path. The work covers both AWS and Azure, and introduces AI agents (Bedrock
Agents, Azure OpenAI Assistants) as first-class entry nodes in the IAM attack
graph, bridging MITRE ATLAS and MITRE ATT\&CK for Cloud. This report acts as a
progress update for the first phase of a two-term Master's project.

\section*{Acknowledgements}
This work has been supported by the academic staff of UNSW's School of Computer
Science and Engineering. Supervision and direction have been provided by
\supervisor.

% ============================================================
% Table of Figures (auto-generated)
% ============================================================
\chapter*{Table of Figures}
\addcontentsline{toc}{chapter}{Table of Figures}
\listoffigures

% ============================================================
% Definitions
% ============================================================
\chapter*{Definitions}
\addcontentsline{toc}{chapter}{Definitions}

\begin{description}
  \item[IAM] Identity and Access Management. The cloud-provider subsystem
    that governs which principals may perform which actions on which resources.
  \item[Attack path] A sequence of permissions and resources by which a
    low-privileged principal can reach an effectively higher-privileged state.
  \item[PoC] Proof of Concept. A working implementation that demonstrates an
    attack primitive end-to-end against a controlled lab.
  \item[IMDS] Instance Metadata Service. The local HTTP endpoint
    (\texttt{169.254.169.254}) on AWS EC2 instances that returns instance and
    role information, including temporary credentials.
  \item[MITRE ATT\&CK for Cloud] The adversary-tactics-and-techniques framework
    extended to IaaS environments.
  \item[MITRE ATLAS] The adversarial-threat framework for AI and machine
    learning systems.
\end{description}

% ============================================================
% Table of Contents (auto-generated)
% ============================================================
\tableofcontents

% ============================================================
\chapter{Introduction}

The rise of cloud computing has shifted the security perimeter from the network
to the identity layer. Misconfigured IAM is now a leading cause of cloud
breaches\cite{placeholder_breach_stats}.

Existing IAM analysis tools largely operate at the per-permission level: they
audit whether a principal holds a given action on a given resource. They miss a
distinct class of vulnerability --- the \emph{combination} of permissions that
forms a viable escalation path. PathTriage addresses this gap by modelling cloud
IAM as a directed graph, discovering reachable escalation paths, ranking them by
a defended exploitability metric, and emitting both detection rules and
mitigation policies for each ranked path.

% ============================================================
\chapter{Previous Work}

\section{IAM Analysis Tools}
% BloodHound OpenGraph, PMapper, IAMSpy, ROADtools

\section{Cloud Attack Path Research}
% Rhino Security Labs catalogue, CloudGoat, MITRE ATT&CK Cloud

\section{AI Agent Security and MITRE ATLAS}
% Prompt injection, function-calling tools as IAM escalation, ATLAS framework

% ============================================================
\chapter{Current Progress}

\section{Lab Architecture}
% baseline + scenarios pattern, per-path artefact convention

\section{Verified Attack Paths}
% Paths 1, 2, 3, ... with MITRE mapping and verification log evidence

\section{PathTriage Prototype}
% CLI + AWS enumerator + NetworkX graph; W2 scope vs deferred edges

\section{Engineering Decision Log}
% audit trail methodology; deviation example (Path 3 self-attach)

% ============================================================
\chapter{Future Work}

\section{Remaining T2 Catalogue}
% paths 4-8, N5 AWS arm (Bedrock chain)

\section{Term 3 Scope}
% Azure deepening, N5 Azure arm, benchmark, expert study, thesis

\section{Risks and Mitigations}
% scope risk, AI service access, vendor dependence

% ============================================================
% Bibliography
% ============================================================
\bibliographystyle{IEEEtran}
\bibliography{refs}

\end{document}
